% === II. RELATED WORK ===
\section{Related Work}
\label{sec:related}

\subsection{Computer Vision-Based Rehabilitation}
Automated rehabilitation has evolved from wearable sensors to vision-based approaches. Chen and Yang~\cite{chen2020pose} pioneered webcam-based exercise form correction using OpenPose keypoints with geometric heuristics and DTW classification, but relied on 2D keypoints sensitive to camera perspective. Pi{\~n}ero-Fuentes et al.~\cite{pinero2021deep} applied deep posture classification to health prevention on edge devices, though targeted office ergonomics rather than rehabilitation. Turner et al.~\cite{appiah2024mobile} proposed mobile form correction via skeleton overlays, while Simoes et al.~\cite{simoes2024accuracy} assessed 2D pose estimation for physiotherapy exercise classification. Comprehensive surveys~\cite{lan2023vision} document the trend of combining vision-based pose analysis with natural language feedback. Existing systems share three limitations that \system{} addresses: (i) reliance on 2D or non-metric pose, limiting joint-angle fidelity; (ii) fitness/ergonomics framing without elderly-specific pain, fatigue, and engagement monitoring; and (iii) absence of voice-interactive coaching in Vietnamese.

\subsection{3D Pose Estimation for Rehabilitation}
Accurate joint angle measurement demands 3D pose estimation. MediaPipe BlazePose~\cite{bazarevsky2020blazepose} is widely adopted for its speed but produces non-metric z-coordinates. Recent models~\cite{zhu2023motionbert,mehraban2024motionagformer,sarandi2020metrabs} improve 3D accuracy; whole-body RTMW3D~\cite{jiang2024rtmw} extends coverage to 133 keypoints (body, hands, face) in real-time, enabling comprehensive exercise monitoring.

\subsection{Pain, Fatigue, and Emotion Detection}
PSPI~\cite{prkachin2008structure} enables pain estimation from AUs, enhanced by deep AU detection methods~\cite{shao2021jaa,wang2025graphau}. Fatigue monitoring uses PERCLOS with Eye Aspect Ratio~\cite{dewi2022adjusting}; engagement detection follows Whitehill et al.~\cite{6786307}. Elderly faces present unique challenges---reduced expressiveness, wrinkles, medication-induced masking---necessitating lightweight approaches.

\subsection{LLM-Based Healthcare Assistance}
LLMs have been explored for rehabilitation guidance, with RAG grounding~\cite{lewis2020retrieval} reducing hallucinations by anchoring responses in clinical literature. Current systems lack integration with real-time pose analysis and voice interaction, particularly for elderly users who benefit from auditory guidance.

\subsection{Biomechanics Metrics for Rehabilitation}
Joint angle computation ranges from dot-product to quaternion-based approaches that handle multi-plane rotations without gimbal lock~\cite{Pirek2017}, with ISB-standardized coordinate systems~\cite{Wu2002} for clinical reporting. Movement smoothness metrics include jerk-based measures~\cite{Rohrer2002} and SPARC, which is duration-independent and validated for stroke rehabilitation~\cite{Scholp2021}. DTW variants enable robust motion comparison across varying execution speeds---essential for elderly users.
